% =====================================================================
% Life Rays & Fragment Chains – Global Configuration (config.tex)
% =====================================================================

% Encoding, fonts, and language
\usepackage[utf8]{inputenc}
\usepackage[T1]{fontenc}
\usepackage[english]{babel}
\usepackage{lmodern}

% Page layout and spacing
\usepackage{geometry}
\geometry{a4paper, margin=2.5cm}

\usepackage{setspace}
\usepackage{parskip}

% Mathematics and graphics
\usepackage{amsmath}
\usepackage{amssymb}
\usepackage{amsfonts}
\usepackage{graphicx}
\graphicspath{{../figures/}}


% Section formatting
\usepackage{titlesec}
\titleformat{\section}{\normalfont\Large\bfseries}{\thesection}{1em}{}


% Header, footer, and page style
\usepackage{fancyhdr}
\usepackage{lastpage}


% Colors and hyperlink setup
\usepackage[dvipsnames]{xcolor}
\usepackage[colorlinks=true, linktoc=all, linkcolor=blue]{hyperref}

% Global hyperref configuration
\hypersetup{
    colorlinks=true,
    linkcolor=blue,
    filecolor=magenta,
    urlcolor=cyan,
    citecolor=ForestGreen,
    pdftitle={Life Rays \& Fragment Chains},
    pdfpagemode=FullScreen,
}

\usepackage[all]{hypcap}
\usepackage{tikz}
\usetikzlibrary{arrows.meta, positioning}


% Glossary system
\usepackage[nonumberlist]{glossaries}
\makeglossaries


% Index system
\usepackage{imakeidx}

% Custom index heading used in the Life Ray Theory style
\newcommand{\indexheading}[1]{%
\hypertarget{index:#1}{\centerline{\Large{#1}}\nopagebreak}%
\pdfbookmark[2]{#1}{index:#1}
}%
\makeindex[options=-s LifeRayTheory2.mst]
